% Research draft during the CP08 freeze. Not integrated or registered.
\paragraph{Exploratory proposition: finitely presented tensor coefficient modules.}
Assume all normalized corollary hypotheses of Claim~\ref{C6}, with the given interval $\alpha:\mathbf1\to J$. Suppose the internal tensor is cocontinuous in each variable. Choose a small $k$-linear indexing category $\mathcal A$, not required to contain zero or finite sums; its additive envelope supplies the corollary's generating-category convention. Write $h_a=\mathcal A(-,a)$. Suppose that for every pair $a,b$ the covariant $\mathcal A$-module
\[
 B_{a,b}(c)=(h_b\otimes h_c)(a)
\]
is finitely presented. Then every chosen projective resolving adjunction $F\dashv G$ induces the desired raw equivalence. No pointwise tensor formula is required.
\begin{proof}
Use the finite relative presentation under $\mathbf1$ constructed in Claim~\ref{C29}. It is a filtered diagram $P_i$ with compatible maps $\alpha_i:\mathbf1\to P_i$ and colimit $\alpha:\mathbf1\to J$, where each stage has a presentation
\[
 \bigoplus_{r\in R_i}h_{b_r}\longrightarrow
 \mathbf1\oplus\bigoplus_{t\in T_i}h_{c_t}\longrightarrow P_i\longrightarrow0
\]
with both displayed index sets finite. The construction used only finite support of representable maps, not the finite-Hom hypothesis of that claim.

Fix $a$. The cocontinuous $k$-linear functor $X\mapsto(P_i\otimes X)(a)$ is tensor with its covariant coefficient module, obtained by evaluating on representables. Cocontinuity of the internal tensor in its first variable gives that coefficient as
\[
 K_i^a=\coker\left(
 \bigoplus_{r\in R_i}B_{a,b_r}\longrightarrow
 \mathcal A(a,-)\oplus\bigoplus_{t\in T_i}B_{a,c_t}\right).
\]
The unit term is exactly $\mathcal A(a,-)$, since $\mathbf1\otimes X=X$. Finite colimits of finitely presented modules are finitely presented, so $K_i^a$ is finitely presented. Its quotient by the image of $\mathcal A(a,-)$ is finitely presented for the same reason.

For every right $\mathcal A$-module $X$, the map $\alpha_i\otimes X$ is monic, because its composite to $J\otimes X$ is $\alpha\otimes X$. Hence the coefficient map
\[
 \mathcal A(a,-)\longrightarrow K_i^a
\]
is pure: its tensor product with every right module is monic. Taking right representables first also verifies that the coefficient map itself is monic. Its finitely presented cokernel therefore splits the pure exact sequence in the category of covariant $\mathcal A$-modules.

For clarity, the finite-equation proof of this splitting works for a small linear category exactly as it does for a ring. Present the cokernel by finitely many generators, each assigned an object of $\mathcal A$, and finitely many relations whose matrix entries are morphisms of $\mathcal A$. Lift the generators to $K_i^a$; the finite relation vector then lies in $\mathcal A(a,-)$. Tensor-purity with the right module presented by the transpose finite matrix supplies a solution in this submodule. Subtracting that solution makes the lifts satisfy the relations, giving a section. Thus no coherence or Noetherian assumption is hidden in the splitting assertion.

Tensoring this coefficient retraction with $X$ gives a splitting of
\[
 X(a)\longrightarrow(P_i\otimes X)(a)
\]
natural in $X$. The splittings need not be natural in $a$ or in $i$. Each $F_n$ is a retract of a sum of evaluations, so $F(\alpha_i\otimes X)$ is degreewise split monic.

For an acyclic nonnegative projective complex $K$ put $Q=GK$. Relative injectivity of $K$ for degreewise split monomorphisms extends the counit across $F(\alpha_i\otimes Q)$. Adjunction gives a retraction $r_i:P_i\otimes Q\to Q$. Cocontinuity gives
\[
 \operatorname{colim}_iF(P_i\otimes Q)=F(J\otimes Q)\simeq0.
\]
Any homology class of $FQ$ maps to a stage and dies at a later stage; applying that later retraction proves it was zero. Claim~\ref{C10} and canonical normalization now give the conclusion by Claim~\ref{C3}.
\end{proof}

\paragraph{Checks of the criterion, not an automatic hypothesis.}
For a coherent pointwise tensor on the chosen original indexing objects, $B_{a,b}=\mathcal A(a,b)\otimes_k\mathcal A(a,-)$, so finitely generated projective Hom modules make these coefficient modules finitely presented. This recovers Claim~\ref{C29}; the pointwise formula is not imposed on sums in the additive envelope. For Day tensor on presheaves over a small monoidal category, $B_{a,b}(c)=\mathcal A(a,b\otimes c)$. If $b$ has a left dual, this is the covariant representable $\mathcal A({}^*b\otimes a,c)$, so the coefficient condition holds for that $b$. Neither closure of finitely presented presheaves under tensor nor ordinary $k$-projectivity of all Hom modules is asserted to imply this coefficient condition.

The rigid Day case is not a nontrivial example of the manuscript hypotheses. If every object of $\mathcal A$ has a dual and $e$ is its monoidal unit, then
\[
 (X\otimes_{\mathrm{Day}}Y)(e)
 \cong\int^{a} X(a)\otimes_kY(a^*).
\]
As $Y$ ranges over right modules, $a\mapsto Y(a^*)$ ranges over all left modules. Universal monicity of $\alpha\otimes Y$ would therefore make $\alpha:h_e\to J$ categorically pure. Tensoring with the covariant weight of $\varphi$ would make $\varphi(\alpha):k\to\varphi J$ monic, contradicting $\varphi J=0$ when $k\ne0$. Thus the rigid specialization merely checks the coefficient formula; it supplies no nonzero-ring interval instance and no strict separation from earlier sufficient criteria.
